% !TeX root = ../queens.tex
\section{The finite verification and its consequences}\label{sec:finite}

We now join the two parts of the proof. By
Lemma~\ref{lem:local-exactness}, the actual branch carries out the
actual next step. The computer checks that
every branch passes the output check and gives a successor satisfying
Condition~\ref{cond:state}. Induction then repeats these two facts for
every column.

\subsection{The starting board}\label{app:seed}

Generate columns $0$ through $29$ directly from the greedy rule.
Their rows $q_0,\dots,q_{29}$ are
\[
\begin{array}{rrrrrrrrrr}
0&2&4&1&3&8&10&12&14&5\\
7&18&6&21&9&24&26&28&30&11\\
13&34&36&38&40&15&17&44&16&47\mathrlap{.}
\end{array}
\]
They give the following state before column $30$:
\begin{center}
\begin{tabular}{ll}
\toprule
Stored field&Value\\
\midrule
$w$&$0$\\
$z$&$-1$\\
$R$&$\varnothing$\\
$D$&$\{1\}$\\
$A$&$\varnothing$\\
$H_{\mathrm{in}}$&\texttt{230121212223}\\
$Q$&\texttt{00322301212}\\
$H_{\mathrm{out}}$&\texttt{300322301212}\\
\bottomrule
\end{tabular}
\end{center}
Here $m=19$, $d=11$, and $\kappa=U(18)=12$, so
$w=30-19-11=0$ and $z=30-19-12=-1$. The input history covers indices $7$--$18$, the
queue $19$--$29$, and the output history $18$--$29$. The verifiers
compute these data from the board rather than accepting the table.
They also check that every lower queen in this prefix satisfies
$|d_j-j|\le4$, and that $\sigma_1\dots\sigma_{29}$ follows the
history graph.

Column $30$ is chosen so that the input history has positive
indices and $\kappa\ge12$, as Lemma~\ref{lem:local-exactness}
requires. On the board $m$ never decreases, so $U(m-1)\ge12$
remains true in every later column. No periodicity of the board is
assumed: Section~\ref{sec:identification} follows the actual queens.

% Layout only: start this subsection, with its algorithm, on a new page.
\newpage
\subsection{Constructing the history graph}\label{sec:finite-statement}

Algorithm~\ref{alg:construction} builds the history graph by running
the calculation of Section~\ref{app:output}, with one difference: a computed symbol that fails the
output check adds its edge and destination vertex to the graph
instead. The algorithm starts from the eighteen windows of
$\sigma_1\dots\sigma_{29}$, those ending at indices $12,\dots,29$, and
their seventeen edges. A successor violating
Condition~\ref{cond:state} ends the construction with failure.

\begin{algorithm}[H]
\caption{Constructing the history graph}
\label{alg:construction}
\renewcommand{\algorithmicensure}{\textbf{Output:}}
\begin{algorithmic}[1]
\Ensure The history graph, or failure.
\State $G\gets$ the windows of $\sigma_1\dots\sigma_{29}$ and the edges between consecutive windows.
\State $\mathcal S\gets\{$the state before column $30\}$.
\Repeat\Comment{one exploration}
  \State Mark every state in $\mathcal S$ as unexamined.
  \While{some state $S\in\mathcal S$ is unexamined}
    \State Mark $S$ as examined.
    \State Run Algorithm~\ref{alg:local-successors} on every branch from $S$, adding each missing output edge to $G$.
    \For{each successor $S'$ found}
      \If{$S'$ violates Condition~\ref{cond:state}}
        \State \Return failure
      \EndIf
      \State Add $S'$ to $\mathcal S$ as unexamined, unless it is already in $\mathcal S$.
    \EndFor
  \EndWhile
\Until{this exploration added no edge to $G$}
\State \Return $G$
\end{algorithmic}
\end{algorithm}

Before column $30$, the queue already holds every symbol the
calculation needs, so the first step makes no request. The queen goes
to $(30,19)$. It is lower, so $u_{30}=0$, and row $30$ contains the
upper queen $(18,30)$, so $b_{30}=1$ and $\sigma_{30}=\mathtt{1}$. The
starting graph has no edge leaving
$H_{\mathrm{out}}=\texttt{300322301212}$, so the output check fails,
and the construction adds the edge
\[
 \texttt{300322301212}\xrightarrow{\;\texttt{1}\;}
 \texttt{003223012121}.
\]

An added edge is a new answer to later requests, so it can create
new branches at states already examined. The construction therefore
repeats the exploration over every state found so far until one adds
no edge. A request with no outgoing edge ends its branch for that
exploration; a later one continues it if an edge has since been
added. In our run, thirteen explorations
add edges and the fourteenth adds none, giving $2092$ vertices and
$2603$ edges.

The result does not depend on the order of exploration. Enlarging
the graph never removes a branch, since a missing edge only ends
one. Hence every edge that the construction adds belongs to every
graph that contains the starting windows and has an edge for each
output that it permits. The final exploration shows that the
constructed graph has this property, so it is the smallest such
graph.

The graph contains many words that never occur in $\sigma$. Because
the state omits most of the board, its records and the symbols it
reads can combine in ways that the board never produces, and the
outputs of all these combinations must be included as well. The
twelve-symbol windows serve two purposes: they retain the
column bits that the attack tests need, and they restrict which
symbols can be read together. With windows of length $4$ through
$11$, the construction reaches a successor violating
Condition~\ref{cond:state} (Appendix~\ref{app:construction}).

% Layout only: start this subsection, with its algorithm, on a new page.
\newpage
\subsection{The exhaustive check}\label{sec:closure-check}

This section supplies what Lemma~\ref{lem:local-exactness} leaves open:
that the output check passes and that the successor satisfies
Condition~\ref{cond:state}. The actual branch cannot be identified in
advance, so the check covers every branch of every state reached.
We now hold the completed history graph fixed and check it with
Algorithm~\ref{alg:check}, adding no edges. The states reached and
their successors form the state graph (Definition~\ref{def:state-graph}).

\begin{algorithm}[H]
\caption{The exhaustive check}
\label{alg:check}
\renewcommand{\algorithmicensure}{\textbf{Output:}}
\begin{algorithmic}[1]
\Ensure The state graph, or failure.
\State Check that every edge of the history graph ends at a vertex, and check the starting board (Section~\ref{app:seed}).
\State $\mathcal S\gets\{$the state before column $30\}$, unexamined; $\mathcal E\gets\varnothing$.
\While{some state $S\in\mathcal S$ is unexamined}
  \State Mark $S$ as examined.
  \State Run Algorithm~\ref{alg:local-successors} on every branch from $S$.\Comment{fails if an output check fails}
  \For{each successor $S'$ found}
    \If{$S'$ violates Condition~\ref{cond:state}}
      \State \Return failure
    \EndIf
    \State Add $S'$ to $\mathcal S$ as unexamined, unless it is already in $\mathcal S$.
    \State Add the edge $S\to S'$ to $\mathcal E$.
  \EndFor
\EndWhile
\State \Return $(\mathcal S,\mathcal E)$
\end{algorithmic}
\end{algorithm}

A state already examined need not be examined again, since its
branches depend only on its eight fields and the fixed graph. The
check has no depth or state-count limit; it ends when no unexamined
state remains. If an output check fails or a successor violates
Condition~\ref{cond:state}, the whole check fails with an error; no
branch is ever dropped to make it pass. Branches that stop at a
vertex without outgoing edges are counted; with the completed graph
there are none.

\begin{proposition}[Finite verification]\label{lem:computation}
The history graph passes the checks above. With this graph fixed,
the calculation from the initial state reaches $7014$ states and
$8327$ directed state-graph edges, and no check fails. In
particular, the initial state satisfies Condition~\ref{cond:state},
and for every reached state, every branch of the calculation passes
the output check and gives a successor satisfying
Condition~\ref{cond:state}.
\end{proposition}
\begin{proof}[Verification]
The tuple implementation explores all these branches and terminates
with no unexamined state or failed check. A second, independent
implementation uses packed bitmasks and performs the same checks. A
comparison program converts both to a common encoding and checks
that the complete state sets and every successor set are equal. The
code and commands are given in Appendix~\ref{app:artifacts}.
\end{proof}

The number $8327$ counts ordered pairs of states, not branches:
several branches can give the same successor, and the row search can
branch again after the queen has been chosen.

Proposition~\ref{lem:computation} concerns the calculation on
states. It remains to show that the actual board stays among these
states; this is the induction below.

% Layout only: start this subsection on a new page.
\newpage
\subsection{Following the actual queens forever}\label{sec:identification}

We keep track of the actual state and the actual queen word
together. Knowing that the earlier word follows the history graph
makes its symbols available as answers to requests. Knowing that the
state is among those checked guarantees that the branch using those
answers has an allowed output and a successor satisfying the bounds.
These are precisely the facts needed to repeat the argument.

\begin{proof}[Proof of Lemma~\ref{lem:estimates}]
We follow a single sequence of states: start with the state before
column $30$ in Section~\ref{app:seed}, and at each column take the
successor given by the actual branch.
Its fields $w,z,R,D,A$ are the actual records, and its words are
actual segments, although the length of $Q$ depends on the requests
made so far. Immediately before column $n\ge30$, we maintain three
assertions about this state:
\begin{enumerate}
\item It is a reached state of the checked state graph, and
satisfies Condition~\ref{cond:state}.
\item The stored word segments are the actual segments at their
stated positions, all with indices below $n$.
\item The actual word through $\sigma_{n-1}$ follows the history
graph.
\end{enumerate}
All three hold at $n=30$ by the direct starting checks.

Assume that they hold before column $n$. Since the least unused row
has not decreased, $U(m-1)\ge12$, so Lemma~\ref{lem:local-exactness}
applies. The actual branch never stops, chooses the actual next queen, and produces its actual symbol
$\sigma_n$. Proposition~\ref{lem:computation} says that this symbol
passes the output check, that is, it labels an edge from the current
$H_{\mathrm{out}}$. Thus the actual word follows the graph through
one more symbol.

The same lemma identifies the successor with the records of the
actual board before column $n+1$, all of whose stored symbols have
indices below $n+1$. Proposition~\ref{lem:computation} places this
successor in the checked collection and supplies
Condition~\ref{cond:state} again. This proves the three assertions
for the next column.

It remains to extract the discrepancy. If the queen just placed is
the $j$th lower queen, with offset $r$, then \eqref{eq:local-rank}
gives
\[
 d_j-j=w-r-|D|.
\]
Condition~\ref{cond:state} gives $0\le w-r\le4$ and $0\le|D|\le4$, so
$|d_j-j|\le4$. The lower queens in columns less than $30$ satisfy the
same bound by the direct prefix check. Hence the bound holds for
every lower queen.
\end{proof}

\begin{corollary}[Diagonal coverage]\label{cor:diagonal-coverage}
Every diagonal $y-x=h$, $h\in\mathbb Z$, contains exactly one
queen. Equivalently, $n\mapsto q_n-n$ is a bijection from $\N$
to $\mathbb Z$. In particular, the lower-diagonal magnitudes
$(d_j)_{j\ge1}$ form a permutation of the positive integers.
\end{corollary}
\begin{proof}
Immediately before the $j$th lower queen is placed in a column
$n\ge30$, the induction gives $D\subseteq[1\dts4]$, and
\eqref{eq:local-rank} gives
\[
 d=j-|D|\ge j-4.
\]
There are infinitely many lower queens, as shown after
\eqref{eq:counts}. The least unused lower-diagonal magnitude
therefore tends to infinity, so every positive magnitude is
eventually used.

Proposition~\ref{prop:stability}, with $C=4$, gives
$U(n)\to\infty$. Lemma~\ref{lem:upper} then puts a queen on
every positive upper diagonal. The origin occupies the main
diagonal, and the nonattacking condition gives uniqueness on
each diagonal.
\end{proof}

\begin{remark}[Zero values on diagonals]
In game terms, every diagonal $y-x=h$ contains exactly one
zero of $G$. This proves the zero-value case of
\cite[Conjecture~24]{DSS}, which asks whether every such diagonal
contains every nonnegative Sprague--Grundy value exactly once.
\end{remark}

\subsection{Consequences for related OEIS sequences}\label{sec:oeis-consequences}

Several OEIS entries~\cite{OEIS} record statistics of the same board
and ask questions about them. Some follow directly from the results
above. Proposition~\ref{prop:stability} proves the upper-queen slope
limit stated in \href{https://oeis.org/A275884}{A275884}. In terms of
our counts, \href{https://oeis.org/A275890}{A275890},
\href{https://oeis.org/A275891}{A275891}, and
\href{https://oeis.org/A275892}{A275892} are $1+U(N-1)$, $L(N-1)$, and
$1+U(N-1)-L(N-1)$, so the same proposition gives $N/\phi+O(1)$,
$N/\phi^2+O(1)$, and $N/\phi^3+O(1)$ for them. By
Corollary~\ref{cor:diagonal-coverage}, the signed diagonal sequences
\href{https://oeis.org/A065185}{A065185} and
\href{https://oeis.org/A276325}{A276325} each enumerate $\mathbb Z$
exactly once.

Other questions concern the pattern of upper and lower columns. They
can be settled from finite graphs derived from the verification: a
graph gives an upper bound for every column after the start, and
direct greedy witnesses show that each bound is attained.

\begin{corollary}[Column runs and gaps]\label{cor:oeis-runs}
The following are the exact sets of values of the indicated sequences.
Runs are maximal; the upper-column sequences include the origin, as in
the OEIS definitions.
\begin{center}
\begin{tabular}{lll}
\toprule
Sequence&Statistic&Set of values\\
\midrule
\href{https://oeis.org/A275885}{A275885}&Lower-column run lengths&$\{1,2,3\}$\\
\href{https://oeis.org/A275886}{A275886}&Upper-column run lengths&$\{1,\ldots,5\}$\\
\href{https://oeis.org/A275887}{A275887}&Run lengths of equal terms in \href{https://oeis.org/A275885}{A275885}&$\{1,\ldots,9,11\}$\\
\href{https://oeis.org/A275888}{A275888}&Gaps between upper columns&$\{1,2,3,4\}$\\
\href{https://oeis.org/A275889}{A275889}&Gaps between lower columns&$\{1,\ldots,6\}$\\
\bottomrule
\end{tabular}
\end{center}
In particular, a $4$ never occurs in \href{https://oeis.org/A275885}{A275885}, and a $10$ never occurs
in \href{https://oeis.org/A275887}{A275887}.
\end{corollary}
\begin{proof}
We first bound the runs and gaps. By the induction of
Section~\ref{sec:identification}, every twelve consecutive symbols of
$\sigma$ form a vertex of the history graph, and no vertex contains
four consecutive symbols with column bit $0$ or six with column
bit~$1$. Hence, among columns $1,2,\dots$, there are never four
consecutive lower columns or six consecutive upper columns; a direct
check covers the origin, which the upper-column sequences count as
upper. So lower runs have length at most three and upper runs at most
five. A gap between consecutive upper columns is one more than the
lower run between them, and similarly for lower columns, so the gaps
between upper columns are at most four and those between lower
columns at most six.

For \href{https://oeis.org/A275887}{A275887}, write $r$ for the run
sequence \href{https://oeis.org/A275885}{A275885} and $g$ for the gap
sequence \href{https://oeis.org/A275888}{A275888}. A gap of $k>1$
between upper columns encloses a lower run of length $k-1$, so $r$ is
obtained from $g$ by deleting its $1$'s and replacing each remaining
term $k$ by $k-1$. The twelve-symbol graph is too coarse to exclude
ten consecutive equal terms in $r$, so we use a refinement.
Appendix~\ref{app:oeis-checks} repeats the verification with
forty-symbol histories and derives from it a finite graph with edges
labeled $1,2,3$, in which the terms of $r$ after column $80$ are the
labels of a walk. For each $c$, the edges labeled $c$ form an acyclic
subgraph, so a maximal run of $c$'s, bounded on both sides by
different labels, has only finitely many possible lengths.
Enumerating these paths gives
\[
 \mathcal L_1=\{1,\ldots,9,11\},\qquad
 \mathcal L_2=\{1,\ldots,6\},\qquad
 \mathcal L_3=\{1\}.
\]
Every run of equal terms in $r$ after column $80$ therefore has
length in $\mathcal L_1\cup\mathcal L_2\cup\mathcal L_3=
\{1,\ldots,9,11\}$, and the earlier runs are checked directly.

Finally, direct greedy witnesses attain every value in the table, so
each set of values is exact.
\end{proof}

These conclusions settle the explicit bound questions in
\href{https://oeis.org/A275885}{A275885} and
\href{https://oeis.org/A275888}{A275888}, and establish the empirical
range recorded in \href{https://oeis.org/A275887}{A275887}.

The same refined graph also settles a catalogue of observations about
$g$. Cut $g$ immediately after each $3$; the pieces are the
\emph{return words} to $3$. They are exactly the $156$ words listed in
the note linked from \href{https://oeis.org/A275888}{A275888}
\cite{OEISWords}: the graph admits no others, and each occurs. Their
lengths range from $2$ to $26$, and exactly five contain a $4$.
Exactly $63$ of them are \emph{faithful}, meaning that every
occurrence is followed by the same return word: for these $63$ the
graph allows only one successor, and for each of the others two
different successors occur. The factors $11111$, $2222$, $33$, $44$,
and $3213$ never occur in $g$. Consecutive $4$'s are at least $71$
terms apart, and every pair at that distance has the same fill,
listed in the note. Appendix~\ref{app:oeis-checks} describes the
finite checks and occurrence witnesses; the complete catalogues are
in the companion repository.

These certificates do not establish the proposed maximum distance
between consecutive $4$'s, limiting word frequencies, or statements
about the positive Sprague--Grundy values.

\subsection{Sharper constants}\label{sec:sharper}
The checked states give better constants than
Theorem~\ref{thm:main}. As in the proof of
Proposition~\ref{prop:stability}, write $\varepsilon(n)=U(n)-n/\phi$.
The following identity replaces the estimate of
Lemma~\ref{lem:counting-recursion} by an exact expression in the
records.

\begin{lemma}[Exact error recursion]\label{lem:exact-recursion}
Consider the actual board before column $n\ge30$, with least unused
row $m$ and records $w$, $R$, and $D$. Let $t$ be the number of lower
rows below $m$. Then $t<n-1$ and
\begin{equation}\label{eq:exact-recursion}
 \phi^2\varepsilon(n-1)=w-|D|-\phi|R|+1+\varepsilon(t).
\end{equation}
\end{lemma}

\begin{proof}
The columns $1,\dots,n-1$ contain $U(n-1)$ upper queens and
$L(n-1)=n-1-U(n-1)$ lower queens. The lower queens use every
magnitude less than $d$ and the $|D|$ magnitudes recorded in $D$, so
$L(n-1)=d-1+|D|$. Since $w=n-m-d$, it follows that
\[
 U(n-1)=m+w-|D|.
\]
Row $m$ is unused. An upper queen in row $m$ would lie in a column
less than $m$, which is already filled, so by Lemma~\ref{lem:surjective}
row $m$ is a lower row. It is therefore the $(t+1)$st lower row, and
Lemma~\ref{lem:complement} gives $m=t+1+U(t)$. The lower rows below
$m$ are all used, necessarily by lower queens, and the lower queens
in rows above $m$ are the $|R|$ recorded in $R$. Hence
$t=L(n-1)-|R|$; in particular $t\le L(n-1)<n-1$, since $q_1=2$ is
upper.

Now eliminate $m$ and $t$. Since $U(t)=t/\phi+\varepsilon(t)$ and
$1+1/\phi=\phi$,
\[
 U(n-1)=w-|D|+1+t+U(t)=w-|D|+1+\phi t+\varepsilon(t).
\]
Substituting $t=n-1-U(n-1)-|R|$ and collecting the multiples of
$U(n-1)$, with $1+\phi=\phi^2$, gives
\[
 \phi^2U(n-1)=\phi(n-1)+w-|D|-\phi|R|+1+\varepsilon(t).
\]
Since $\phi^2U(n-1)=\phi(n-1)+\phi^2\varepsilon(n-1)$, this is
\eqref{eq:exact-recursion}.
\end{proof}

The verification also checks that every state of the state graph
satisfies
\begin{equation}\label{eq:sharper-check}
 -4\le w-|D|-\phi|R|\le1,
\end{equation}
and that $w-r-|D|\le2$ at every lower choice
(Appendix~\ref{app:verification}). With these facts, the recursion
\eqref{eq:exact-recursion} gives the following bounds.

\begin{proposition}[Sharper constants]\label{prop:sharper}
For every $n\ge0$, $-3/\phi<\varepsilon(n)<2/\phi$. Consequently, for
every $n\ge1$,
\[
 \begin{aligned}
 1-\frac{4}{\phi}&<q_n-n\phi<\frac{2}{\phi}&&\text{if }q_n>n,\\[3pt]
 -2-\frac{4}{\phi}&<q_n-\frac{n}{\phi}<4+\frac{1}{\phi}&&\text{if }q_n<n.
 \end{aligned}
\]
\end{proposition}

\begin{proof}
For $0\le n\le29$, the bounds on $\varepsilon(n)$ are checked
directly. Let $n\ge30$, and assume them for all smaller arguments. By
the induction of Section~\ref{sec:identification}, the state before
column $n+1$ belongs to the state graph, so its records satisfy
\eqref{eq:sharper-check}, and Lemma~\ref{lem:exact-recursion} gives
\[
 \phi^2\varepsilon(n)=w-|D|-\phi|R|+1+\varepsilon(t)
\]
for some $t<n$. The induction hypothesis and \eqref{eq:sharper-check}
then give
\[
 \phi^2\varepsilon(n)<1+1+\frac{2}{\phi}=2\phi,\qquad
 \phi^2\varepsilon(n)>-4+1-\frac{3}{\phi}=-3\phi,
\]
and dividing by $\phi^2$ gives the bounds on $\varepsilon(n)$.

If $q_n>n$, then $q_n-n\phi=\varepsilon(n)$, as in the proof of
Proposition~\ref{prop:stability}. Column $n$ adds one upper queen, so
$\varepsilon(n)=\varepsilon(n-1)+1-1/\phi=\varepsilon(n-1)+1/\phi^2$,
which is greater than $1/\phi^2-3/\phi=1-4/\phi$. The upper bound is
the bound on $\varepsilon(n)$ itself.

If column $n$ contains the $j$th lower queen, then
$q_n-n/\phi=\varepsilon(n)+j-d_j$ by \eqref{eq:rank-row}. Here
$\varepsilon(n)=\varepsilon(n-1)-1/\phi$ lies strictly between
$-4/\phi$ and $1/\phi$. Lemma~\ref{lem:estimates} gives $d_j-j\ge-4$.
From column $30$ on, $d_j-j=w-r-|D|\le2$ by \eqref{eq:local-rank},
and the lower queens before column $30$ satisfy the same bound by
direct check. Hence $-4/\phi-2<q_n-n/\phi<1/\phi+4$.
\end{proof}

In the $1$-indexed coordinates $c=n+1$ and $s(c)=q_n+1$ of the remark
after Proposition~\ref{prop:stability}, the proposition gives, after
rounding, for every $c\ge1$,
\[
 s(c)\in
 \Bigl[\frac{c}{\phi}-4.091\dts\frac{c}{\phi}+5\Bigr]
 \cup[c\phi-2.091\dts c\phi+0.619].
\]
Compared with Knuth's ranges
$[c/\phi-3\dts c/\phi+5]\cup[c\phi-2\dts c\phi+1]$, the upper endpoints
are no larger, so the upper halves of those ranges hold for every
$c\ge1$. The lower halves remain open.
